\section{Carry and incidence registers and projective reconstruction of histories}
\label{lib07:armonizacion}
The APP--TRIT--TPK construction fixes an enriched state before
its numerical, incidence and operator realizations. The registers
entering that construction are distinguished by the operation
they allow to be inverted and by the relations they preserve. This section
brings these distinctions together and composes recovery with refinement.

\subsection{Carry, history, incidence and archive}
The carry in Sections~\ref{lib03:residuo} and
\ref{lib03:memoria} is an integer arithmetic datum. Preserving it,
together with depth and the final state, restores the antecedent
through the explicit inverse~\eqref{lib03:inversa-afin}.
The transport history contains the ordered sequence of operations,
sheets, orientations and returns. Within it, two stages with the same phase
may have different accumulated registers. Holonomy acts on
that history and on the frame in which its readings are performed.

The incidence register preserves the boundary relations
used by the exceptional constructions. The ordered object \(K\)
and its reading \(\kappa_{\rm ag}\) enter with the normalizations
and maps of Section~\ref{lib01:identificacion}.
Sufficiency for recovering an operator is expressed there through
the frame responses and the corresponding inversion.
The archive memory of Section~\ref{lib04:archivo}, in turn,
is a sequence of components of a Hilbert space. The telescoping
identity~\eqref{lib04:telescopia} preserves the norm between the
remaining state and the archived components; the inverse
\eqref{lib04:recuperacion-infinita} reconstructs the initial amplitude.

These four functions are composed while preserving their types.
A carry integer, a transport word, an incidence
and an archived amplitude constitute distinct data of the system.
The recovery proofs specify which combination of them
suffices for each reading. Figure~\ref{lib07:fig-registros}
represents this functional hierarchy.

\begin{figure}[htbp]
\centering
\begin{tikzpicture}[
 node distance=9mm,
 box/.style={draw=ink,rounded corners=2pt,fill=white,align=center,
 text width=58mm,minimum height=15mm,font=\small},
 arr/.style={-{Latex[length=2mm]},draw=ink,thick}]
\node[box,text width=125mm] (s) {Enriched APP--TRIT--TPK state\\
 operation, sheet, orientation, boundary and memory};
\node[box] (a) at (-33mm,-30mm)
 {Arithmetic carry\\Remainder, depth and inverse};
\node[box] (h) at (33mm,-30mm)
 {Transport history\\Ordered composition and holonomy};
\node[box] (q) at (-33mm,-62mm)
 {Incidence register\\Boundary and frame readings};
\node[box] (m) at (33mm,-62mm)
 {Recovery archive\\Components and reconstruction};
\node[box,text width=125mm] (r) at (0,-94mm)
 {Recovery compatible with refinement\\
 Typed inverses and preserved relations};
\draw[arr] (s.south -| a.north)--(a);
\draw[arr] (s.south -| h.north)--(h);
\draw[arr] (a)--(q);
\draw[arr] (h)--(m);
\draw[arr] (q.south)--(r.north -| q.south);
\draw[arr] (m.south)--(r.north -| m.south);
\end{tikzpicture}
\caption{Organization of the registers by their role.
The lower arrows indicate composition of recovery
conditions, not identification of the four data types.}
\label{lib07:fig-registros}
\end{figure}

\subsection{Refinement tree and local weights}
Let \(H_k\) be the finite set of admissible prefixes of depth
\(k\), with truncations \(p_{k+1,k}\), of the history system already
constructed. An emission is counted with its multiplicity, and
\(d(\zeta)>0\) denotes the number of emissions from \(\zeta\).
The uniform seed distribution and the local weight
\(1/d(\zeta)\) define \(\mu_k\). The sum of the weights of the
extensions of a prefix equals its weight: the measures are
compatible and determine the cylinder measure \(\mu\).
The propositions on fiber memory and
naturality of the measure are proved below using these explicit definitions.

\subsection{Enriched readers and fiber memory}
\label{sec:ch-memoria-fibra}

The difference between history and value admits an operational criterion. At a
finite level \(H_k\), let \(q_k:H_k\to Z_k\) be a reader and let
\(M_k:H_k\to\mathcal M_k\) be the retained register. The enriched reader
is \(\widehat q_k=(q_k,M_k)\). Memory distinguishes states within the same
reader fiber; its sufficiency is checked on the domain under consideration.
This formulation incorporates into the article itself the complete development of the
fiber information used below.

\begin{proposition}[Recoverability at a finite level]
\label{prop:ch-recuperabilidad-fibra}
An inverse of \(\widehat q_k\) on its image exists if and only if
\(\widehat q_k\) is injective. For a random variable \(Z\) with
values in \(H_k\) and distribution \(\mu_k\), we have
\[
 \mathsf H(Z)=\mathsf H(q_k(Z))+\mathsf H(Z\mid q_k(Z)).
\]
If \(\widehat q_k\) is injective on the support of \(Z\), then
\[
 \mathsf H(Z\mid q_k(Z),M_k(Z))=0.
\]
Here \(\mathsf H(Z)=-\sum_z\Pr(Z=z)\log\Pr(Z=z)\), with
\(0\log0=0\), is finite entropy in nats; conditional entropy is
the average entropy of the conditional distributions.
\end{proposition}

\begin{proof}
An inverse recovers a unique antecedent, which requires injectivity;
conversely, injectivity allows each expressed pair to be assigned its unique
antecedent. Since \(q_k(Z)\) is a deterministic function of \(Z\),
grouping the sum defining \(\mathsf H(Z)\) by fibers of \(q_k\) gives the
chain rule. If the pair \((q_k(Z),M_k(Z))\) distinguishes the support,
each corresponding conditional distribution is concentrated on a single
state and its entropy is zero.
\end{proof}

Not publishing a coordinate and erasing the register that allows it to be recovered are
distinct operations. This result neither identifies entropy with the cardinality
of the continuum nor introduces a thermal interpretation: its domain, reader and
distribution are those just fixed.

\subsection{Equivariance of the measure under tree symmetries}
\label{sec:ch-naturalidad-medida}

\begin{proposition}[Joint transport of the tree and its measure]
\label{prop:ch-naturalidad-medida}
Let \(g=(g_k)_{k\geq0}\) be a collection of bijections \(g_k:H_k\to H_k\)
commuting with truncations, fixing the root and transporting
admissible emissions bijectively while preserving their multiplicities. Then
\[
 d(g_k\zeta)=d(\zeta),\qquad
 (g_k)_*\mu_k=\mu_k,\qquad (g_\infty)_*\mu=\mu.
\]
More generally, if \(\mu_\zeta\) is the law of continuations of a
history \(\zeta\in H_k\), we have
\[
 (g_\infty)_*\mu_\zeta=\mu_{g_k\zeta}.
\]
\end{proposition}

\begin{proof}
The bijection of emissions preserves their number with multiplicities. The
uniform seed distribution is also invariant. Each trajectory
and its image therefore receive the same product of local weights
\(1/d(\zeta)\). Equality of measures follows at each level and on
each cylinder. Coherence with truncations defines \(g_\infty\),
and uniqueness of the measure determined by cylinders proves equality
in the limit. The same argument, starting at \(\zeta\), proves the
assertion about continuations.
\end{proof}

This gives the naturality of the tree and measure: a coherent change
of representation transports the weights without requiring them to be chosen anew.
Preservation of emissions and multiplicities is an essential hypothesis of
the proposition, not a datum added after the construction.

\subsection{Recovery of histories and transport of membership}
\begin{proposition}[Compatibility of inverses and collections]
\label{lib07:prop-familias}
Let \(H_k\) be the preceding levels and \(Y_k=\widehat q_k(H_k)\).
Suppose that \(\widehat q_k:H_k\to Y_k\) is bijective and that there exist
truncations \(s_{k+1,k}\) such that
\[
 s_{k+1,k}\widehat q_{k+1}=\widehat q_k p_{k+1,k}.
\]
Then the finite inverses induce a bijection
\[
 \widehat q_\infty:\varprojlim H_k\longrightarrow\varprojlim Y_k.
\]
If \(H=\varprojlim H_k\), \(Y=\varprojlim Y_k\) and
\(\mathcal T(A)=\widehat q_\infty[A]\), the map
\(\mathcal T:\mathcal P(H)\to\mathcal P(Y)\) is bijective and
preserves unions, intersections, relative complements and membership.
Each restriction \(A\to\mathcal T(A)\) is a bijection.
\end{proposition}
\begin{proof}
Composing the compatibility identity with the inverses gives
\[
 p_{k+1,k}\widehat q_{k+1}^{-1}
   =\widehat q_k^{-1}s_{k+1,k}.
\]
Therefore, a compatible collection \(y=(y_k)\) determines the compatible
collection \(x=(\widehat q_k^{-1}(y_k))\). The two coordinate
maps are inverses. Direct image under a bijection has
as its inverse direct image under the inverse bijection. For every
collection \((A_i)_{i\in I}\) of subsets of \(H\),
\[
 \begin{aligned}
 \mathcal T\Bigl(\bigcup_i A_i\Bigr)&=\bigcup_i\mathcal T(A_i),\\
 \mathcal T\Bigl(\bigcap_i A_i\Bigr)&=\bigcap_i\mathcal T(A_i),\\
 \mathcal T(H\setminus A)&=Y\setminus\mathcal T(A).
 \end{aligned}
\]
The first equality uses the definition of image. For the second and
third, uniqueness of the antecedent allows the membership
conditions to be transferred. Finally,
\(x\in A\) is equivalent to
\(\widehat q_\infty(x)\in\mathcal T(A)\), which also proves
the assertion about restrictions.
\end{proof}

The route-refinement theorem~\ref{lib05:teo-rutas}
provides a realization of this compatibility on the arithmetic
labels and their carries. The result distinguishes recovery
of each state from preservation of operations on collections.
At amplitude level, the unitary lift
in~\ref{lib05:prop-unitaria} also transports the inner product.
At probabilistic level, Proposition
\ref{prop:ch-naturalidad-medida} preserves the weights when
emissions are transported with their multiplicities.

\subsection{Connection with the genealogical closure of the continuum}
Section~\ref{sec:cierre-cardinal-nativo} constructs a genealogical
closure \(\mathfrak G\) and establishes the cardinal identity
\(2^{\aleph_0}=\aleph_1\) over its internal power set. The closure,
generator code and realized history form part of its
definition; its set-theoretic representation is established after the
genealogical construction through the operations and proofs
included in that section.

The present connection incorporates the proofs concerning memory,
naturality and transport of collections used by that relation with
conservation. The proof chain remains differentiated:
the inverses recover states; their compatibility recovers
histories; the bijection transports membership operations;
the genealogical closure in Section~\ref{sec:cierre-cardinal-nativo} fixes the domain
of the cardinal statement. Recovery and the organization
of collections are thus integrated into the double projection with their respective
objects, operations and conclusions.
